\IfFileExists{stacks-project.cls}{%
\documentclass{stacks-project}
}{%
\documentclass{amsart}
}

% For dealing with references we use the comment environment
\usepackage{verbatim}
\newenvironment{reference}{\comment}{\endcomment}
%\newenvironment{reference}{}{}
\newenvironment{slogan}{\comment}{\endcomment}
\newenvironment{history}{\comment}{\endcomment}

% For commutative diagrams we use Xy-pic
\usepackage[all]{xy}

% We use 2cell for 2-commutative diagrams.
\xyoption{2cell}
\UseAllTwocells

% We use multicol for the list of chapters between chapters
\usepackage{multicol}

% This is generally recommended for better output
\usepackage{lmodern}
\usepackage[T1]{fontenc}

% For cross-file-references
\usepackage{xr-hyper}

% Package for hypertext links:
\usepackage{hyperref}

% For any local file, say "hello.tex" you want to link to please
% use \externaldocument[hello-]{hello}
\externaldocument[introduction-]{introduction}
\externaldocument[conventions-]{conventions}
\externaldocument[sets-]{sets}
\externaldocument[categories-]{categories}
\externaldocument[topology-]{topology}
\externaldocument[sheaves-]{sheaves}
\externaldocument[sites-]{sites}
\externaldocument[stacks-]{stacks}
\externaldocument[fields-]{fields}
\externaldocument[algebra-]{algebra}
\externaldocument[brauer-]{brauer}
\externaldocument[homology-]{homology}
\externaldocument[derived-]{derived}
\externaldocument[simplicial-]{simplicial}
\externaldocument[more-algebra-]{more-algebra}
\externaldocument[smoothing-]{smoothing}
\externaldocument[modules-]{modules}
\externaldocument[sites-modules-]{sites-modules}
\externaldocument[injectives-]{injectives}
\externaldocument[cohomology-]{cohomology}
\externaldocument[sites-cohomology-]{sites-cohomology}
\externaldocument[dga-]{dga}
\externaldocument[dpa-]{dpa}
\externaldocument[sdga-]{sdga}
\externaldocument[hypercovering-]{hypercovering}
\externaldocument[schemes-]{schemes}
\externaldocument[constructions-]{constructions}
\externaldocument[properties-]{properties}
\externaldocument[morphisms-]{morphisms}
\externaldocument[coherent-]{coherent}
\externaldocument[divisors-]{divisors}
\externaldocument[limits-]{limits}
\externaldocument[varieties-]{varieties}
\externaldocument[topologies-]{topologies}
\externaldocument[descent-]{descent}
\externaldocument[perfect-]{perfect}
\externaldocument[more-morphisms-]{more-morphisms}
\externaldocument[flat-]{flat}
\externaldocument[groupoids-]{groupoids}
\externaldocument[more-groupoids-]{more-groupoids}
\externaldocument[etale-]{etale}
\externaldocument[chow-]{chow}
\externaldocument[intersection-]{intersection}
\externaldocument[pic-]{pic}
\externaldocument[weil-]{weil}
\externaldocument[adequate-]{adequate}
\externaldocument[dualizing-]{dualizing}
\externaldocument[duality-]{duality}
\externaldocument[discriminant-]{discriminant}
\externaldocument[derham-]{derham}
\externaldocument[local-cohomology-]{local-cohomology}
\externaldocument[algebraization-]{algebraization}
\externaldocument[curves-]{curves}
\externaldocument[resolve-]{resolve}
\externaldocument[models-]{models}
\externaldocument[functors-]{functors}
\externaldocument[equiv-]{equiv}
\externaldocument[pione-]{pione}
\externaldocument[etale-cohomology-]{etale-cohomology}
\externaldocument[proetale-]{proetale}
\externaldocument[relative-cycles-]{relative-cycles}
\externaldocument[more-etale-]{more-etale}
\externaldocument[trace-]{trace}
\externaldocument[crystalline-]{crystalline}
\externaldocument[spaces-]{spaces}
\externaldocument[spaces-properties-]{spaces-properties}
\externaldocument[spaces-morphisms-]{spaces-morphisms}
\externaldocument[decent-spaces-]{decent-spaces}
\externaldocument[spaces-cohomology-]{spaces-cohomology}
\externaldocument[spaces-limits-]{spaces-limits}
\externaldocument[spaces-divisors-]{spaces-divisors}
\externaldocument[spaces-over-fields-]{spaces-over-fields}
\externaldocument[spaces-topologies-]{spaces-topologies}
\externaldocument[spaces-descent-]{spaces-descent}
\externaldocument[spaces-perfect-]{spaces-perfect}
\externaldocument[spaces-more-morphisms-]{spaces-more-morphisms}
\externaldocument[spaces-flat-]{spaces-flat}
\externaldocument[spaces-groupoids-]{spaces-groupoids}
\externaldocument[spaces-more-groupoids-]{spaces-more-groupoids}
\externaldocument[bootstrap-]{bootstrap}
\externaldocument[spaces-pushouts-]{spaces-pushouts}
\externaldocument[spaces-chow-]{spaces-chow}
\externaldocument[groupoids-quotients-]{groupoids-quotients}
\externaldocument[spaces-more-cohomology-]{spaces-more-cohomology}
\externaldocument[spaces-simplicial-]{spaces-simplicial}
\externaldocument[spaces-duality-]{spaces-duality}
\externaldocument[formal-spaces-]{formal-spaces}
\externaldocument[restricted-]{restricted}
\externaldocument[spaces-resolve-]{spaces-resolve}
\externaldocument[formal-defos-]{formal-defos}
\externaldocument[defos-]{defos}
\externaldocument[cotangent-]{cotangent}
\externaldocument[examples-defos-]{examples-defos}
\externaldocument[algebraic-]{algebraic}
\externaldocument[examples-stacks-]{examples-stacks}
\externaldocument[stacks-sheaves-]{stacks-sheaves}
\externaldocument[criteria-]{criteria}
\externaldocument[artin-]{artin}
\externaldocument[quot-]{quot}
\externaldocument[stacks-properties-]{stacks-properties}
\externaldocument[stacks-morphisms-]{stacks-morphisms}
\externaldocument[stacks-limits-]{stacks-limits}
\externaldocument[stacks-cohomology-]{stacks-cohomology}
\externaldocument[stacks-perfect-]{stacks-perfect}
\externaldocument[stacks-introduction-]{stacks-introduction}
\externaldocument[stacks-more-morphisms-]{stacks-more-morphisms}
\externaldocument[stacks-geometry-]{stacks-geometry}
\externaldocument[moduli-]{moduli}
\externaldocument[moduli-curves-]{moduli-curves}
\externaldocument[examples-]{examples}
\externaldocument[exercises-]{exercises}
\externaldocument[guide-]{guide}
\externaldocument[desirables-]{desirables}
\externaldocument[coding-]{coding}
\externaldocument[obsolete-]{obsolete}
\externaldocument[fdl-]{fdl}
\externaldocument[index-]{index}

% Theorem environments.
%
\theoremstyle{plain}
\newtheorem{theorem}[subsection]{Theorem}
\newtheorem{proposition}[subsection]{Proposition}
\newtheorem{lemma}[subsection]{Lemma}

\theoremstyle{definition}
\newtheorem{definition}[subsection]{Definition}
\newtheorem{example}[subsection]{Example}
\newtheorem{exercise}[subsection]{Exercise}
\newtheorem{situation}[subsection]{Situation}

\theoremstyle{remark}
\newtheorem{remark}[subsection]{Remark}
\newtheorem{remarks}[subsection]{Remarks}

\numberwithin{equation}{subsection}

% Macros
%
\def\lim{\mathop{\mathrm{lim}}\nolimits}
\def\colim{\mathop{\mathrm{colim}}\nolimits}
\def\Spec{\mathop{\mathrm{Spec}}}
\def\Hom{\mathop{\mathrm{Hom}}\nolimits}
\def\Ext{\mathop{\mathrm{Ext}}\nolimits}
\def\SheafHom{\mathop{\mathcal{H}\!\mathit{om}}\nolimits}
\def\SheafExt{\mathop{\mathcal{E}\!\mathit{xt}}\nolimits}
\def\Sch{\mathit{Sch}}
\def\Mor{\mathop{\mathrm{Mor}}\nolimits}
\def\Ob{\mathop{\mathrm{Ob}}\nolimits}
\def\Sh{\mathop{\mathit{Sh}}\nolimits}
\def\NL{\mathop{N\!L}\nolimits}
\def\CH{\mathop{\mathrm{CH}}\nolimits}
\def\proetale{{pro\text{-}\acute{e}tale}}
\def\etale{{\acute{e}tale}}
\def\QCoh{\mathit{QCoh}}
\def\Ker{\mathop{\mathrm{Ker}}}
\def\Im{\mathop{\mathrm{Im}}}
\def\Coker{\mathop{\mathrm{Coker}}}
\def\Coim{\mathop{\mathrm{Coim}}}

% Boxtimes
%
\DeclareMathSymbol{\boxtimes}{\mathbin}{AMSa}{"02}

%
% Macros for moduli stacks/spaces
%
\def\QCohstack{\mathcal{QC}\!\mathit{oh}}
\def\Cohstack{\mathcal{C}\!\mathit{oh}}
\def\Spacesstack{\mathcal{S}\!\mathit{paces}}
\def\Quotfunctor{\mathrm{Quot}}
\def\Hilbfunctor{\mathrm{Hilb}}
\def\Curvesstack{\mathcal{C}\!\mathit{urves}}
\def\Polarizedstack{\mathcal{P}\!\mathit{olarized}}
\def\Complexesstack{\mathcal{C}\!\mathit{omplexes}}
% \Pic is the operator that assigns to X its picard group, usage \Pic(X)
% \Picardstack_{X/B} denotes the Picard stack of X over B
% \Picardfunctor_{X/B} denotes the Picard functor of X over B
\def\Pic{\mathop{\mathrm{Pic}}\nolimits}
\def\Picardstack{\mathcal{P}\!\mathit{ic}}
\def\Picardfunctor{\mathrm{Pic}}
\def\Deformationcategory{\mathcal{D}\!\mathit{ef}}
\usepackage{cleveref}
\crefname{originalremark}{original remark}{original remarks}
\Crefname{originalremark}{Original remark}{Original remarks}
\makeatletter
\newlabel{situation-relative-descent}{{0DI6}{source}{}{sourcepointer-1}{}}
\newlabel{situation-relative-descent@cref}{{[originalremark][0][]0DI6}{[1][0][]source}{}{}{}}
\makeatother
\crefname{originallemma}{original lemma}{original lemmas}
\Crefname{originallemma}{Original lemma}{Original lemmas}
\makeatletter
\newlabel{cohomology-lemma-mayer-vietoris-hom}{{08BW}{source}{}{sourcepointer-2}{}}
\newlabel{cohomology-lemma-mayer-vietoris-hom@cref}{{[originallemma][0][]08BW}{[1][0][]source}{}{}{}}
\makeatother
\crefname{originallemma}{original lemma}{original lemmas}
\Crefname{originallemma}{Original lemma}{Original lemmas}
\makeatletter
\newlabel{coherent-lemma-induction-principle}{{08DR}{source}{}{sourcepointer-3}{}}
\newlabel{coherent-lemma-induction-principle@cref}{{[originallemma][0][]08DR}{[1][0][]source}{}{}{}}
\makeatother
\crefname{originallemma}{original lemma}{original lemmas}
\Crefname{originallemma}{Original lemma}{Original lemmas}
\makeatletter
\newlabel{more-algebra-lemma-colimit-relatively-perfect}{{0DHX}{source}{}{sourcepointer-4}{}}
\newlabel{more-algebra-lemma-colimit-relatively-perfect@cref}{{[originallemma][0][]0DHX}{[1][0][]source}{}{}{}}
\makeatother
\crefname{originallemma}{original lemma}{original lemmas}
\Crefname{originallemma}{Original lemma}{Original lemmas}
\makeatletter
\newlabel{lemma-pseudo-coherent}{{08E5}{source}{}{sourcepointer-5}{}}
\newlabel{lemma-pseudo-coherent@cref}{{[originallemma][0][]08E5}{[1][0][]source}{}{}{}}
\makeatother
\crefname{originallemma}{original lemma}{original lemmas}
\Crefname{originallemma}{Original lemma}{Original lemmas}
\makeatletter
\newlabel{lemma-affine-compare-bounded}{{06Z0}{source}{}{sourcepointer-6}{}}
\newlabel{lemma-affine-compare-bounded@cref}{{[originallemma][0][]06Z0}{[1][0][]source}{}{}{}}
\makeatother
\crefname{originallemma}{original lemma}{original lemmas}
\Crefname{originallemma}{Original lemma}{Original lemmas}
\makeatletter
\newlabel{lemma-pseudo-coherent-affine}{{08E7}{source}{}{sourcepointer-7}{}}
\newlabel{lemma-pseudo-coherent-affine@cref}{{[originallemma][0][]08E7}{[1][0][]source}{}{}{}}
\makeatother
\crefname{originallemma}{original lemma}{original lemmas}
\Crefname{originallemma}{Original lemma}{Original lemmas}
\makeatletter
\newlabel{lemma-tor-dimension-rel-affine}{{0DHY}{source}{}{sourcepointer-8}{}}
\newlabel{lemma-tor-dimension-rel-affine@cref}{{[originallemma][0][]0DHY}{[1][0][]source}{}{}{}}
\makeatother
\title{Descent of Hom groups out of a pseudo-coherent object into a relatively finite tor dimension object along a filtered limit of schemes}
\author{The Stacks Project Authors}
\date{stacks-project:perfect.tex:local-snapshot-20261012; GNU FDL 1.2 or later}
\begin{document}
\maketitle
\setcounter{subsection}{1}
\noindent\textit{Editorial note.} All delivered mathematical text is the author's own, copied from the pinned Stacks Project chapter file perfect.tex; no statement or proof step was rewritten, reordered or removed. The counted span is the author environment and its own proof environment. References inside the span are handed to the assembler as navigation stubs that carry the canonical tag of the cited tag only; whether each of them is load-bearing is recorded on the stub itself.\par\medskip
\noindent\textit{Allowed standard prerequisites (editorial).}\begin{itemize}
\item Mayer-Vietoris for Hom in the derived category: Hom out of a compact-free source can be computed across unions of opens, so a local isomorphism statement propagates from U\_0, V\_0 and U\_0 cap V\_0 to U\_0 cup V\_0. invoked at proof lines 9082-9085. Provides the gluing step of the local condition P.
\item The induction principle for quasi-compact opens, applied first to opens of S\_0 and then to quasi-compact quasi-separated opens of X\_0. invoked at proof lines 9090-9098. Reduces the assertion to the affine-affine case.
\item In the affine case the derived-category Hom statement translates into module Hom computations, with K\_0 pseudo-coherent and L\_0 of finite tor dimension over the ring, by the pseudo-coherence, affine-compare, pseudo-coherent-affine, and tor-dimension lemmas. invoked at proof lines 9101-9112. Sets up the hypotheses of the More on Algebra colimit lemma for relatively perfect systems.
\end{itemize}\medskip
\ifdefined\hypertarget\hypertarget{sourcepointer-1}{}\fi
\noindent\textit{Source reference.} Original remark 0DI6: sources/stacks-chow/tags.txt, source lines 16931--16931. Cited at proof line 9058 as the standing situation of the lemma; the gluing and reduction arguments are inside the counted span.\par\medskip
\ifdefined\hypertarget\hypertarget{sourcepointer-2}{}\fi
\noindent\textit{Source reference.} Original lemma 08BW: sources/stacks-chow/tags.txt, source lines 10581--10581. Cited at proof lines 9084-9085 to glue the condition P over unions; the gluing argument is inside the counted span.\par\medskip
\ifdefined\hypertarget\hypertarget{sourcepointer-3}{}\fi
\noindent\textit{Source reference.} Original lemma 08DR: sources/stacks-chow/tags.txt, source lines 10646--10646. Cited at proof lines 9091-9092 and used again at proof lines 9096-9098 to reduce to the affine-affine case; both reductions are inside the counted span.\par\medskip
\ifdefined\hypertarget\hypertarget{sourcepointer-4}{}\fi
\noindent\textit{Source reference.} Original lemma 0DHX: sources/stacks-chow/tags.txt, source lines 16922--16922. Cited at proof lines 9102 and 9114 to conclude in the affine case; the translation into module terms is inside the counted span.\par\medskip
\ifdefined\hypertarget\hypertarget{sourcepointer-5}{}\fi
\noindent\textit{Source reference.} Original lemma 08E5: sources/stacks-perfect/pinned-original-source.tex, source lines 2422--2422. Cited at proof line 9103 to translate the statement; the module-level translation is inside the counted span.\par\medskip
\ifdefined\hypertarget\hypertarget{sourcepointer-6}{}\fi
\noindent\textit{Source reference.} Original lemma 06Z0: sources/stacks-perfect/pinned-original-source.tex, source lines 454--454. Cited at proof lines 9103-9104 to translate Homs into module Homs; the translation is inside the counted span.\par\medskip
\ifdefined\hypertarget\hypertarget{sourcepointer-7}{}\fi
\noindent\textit{Source reference.} Original lemma 08E7: sources/stacks-perfect/pinned-original-source.tex, source lines 2440--2440. Cited at proof lines 9107-9109 for pseudo-coherence of the module complex of K\_0; the affine case assembly is inside the counted span.\par\medskip
\ifdefined\hypertarget\hypertarget{sourcepointer-8}{}\fi
\noindent\textit{Source reference.} Original lemma 0DHY: sources/stacks-perfect/pinned-original-source.tex, source lines 2552--2552. Cited at proof lines 9110-9112 for the finite tor dimension of the module complex of L\_0; the affine case assembly is inside the counted span.\par\medskip
\begin{lemma}
\label{lemma-relative-descend-homomorphisms}
In Situation \ref{situation-relative-descent}.
Let $K_0$ and $L_0$ be objects of $D(\mathcal{O}_{X_0})$.
Set $K_i = Lf_{i0}^*K_0$ and $L_i = Lf_{i0}^*L_0$ for $i \geq 0$
and set $K = Lf_0^*K_0$ and $L = Lf_0^*L_0$. Then the map
$$
\colim_{i \geq 0} \Hom_{D(\mathcal{O}_{X_i})}(K_i, L_i)
\longrightarrow
\Hom_{D(\mathcal{O}_X)}(K, L)
$$
is an isomorphism if $K_0$ is pseudo-coherent and
$L_0 \in D_\QCoh(\mathcal{O}_{X_0})$ has (locally)
finite tor dimension as an object of
$D((X_0 \to S_0)^{-1}\mathcal{O}_{S_0})$
\end{lemma}

\begin{proof}
For every quasi-compact open $U_0 \subset X_0$ consider the
condition $P$ that
$$
\colim_{i \geq 0} \Hom_{D(\mathcal{O}_{U_i})}(K_i|_{U_i}, L_i|_{U_i})
\longrightarrow
\Hom_{D(\mathcal{O}_U)}(K|_U, L|_U)
$$
is an isomorphism where $U = f_0^{-1}(U_0)$ and $U_i = f_{i0}^{-1}(U_0)$.
If $P$ holds for $U_0$, $V_0$ and $U_0 \cap V_0$, then it holds
for $U_0 \cup V_0$ by Mayer-Vietoris
for hom in the derived category, see
Cohomology, Lemma \ref{cohomology-lemma-mayer-vietoris-hom}.

\medskip\noindent
Denote $\pi_0 : X_0 \to S_0$ the given morphism.
Then we can first consider $U_0 = \pi_0^{-1}(W_0)$ with
$W_0 \subset S_0$ quasi-compact open. By the induction principle of
Cohomology of Schemes, Lemma \ref{coherent-lemma-induction-principle}
applied to quasi-compact opens of $S_0$
and the remark above, we find that it is enough to prove
$P$ for $U_0 = \pi_0^{-1}(W_0)$ with $W_0$ affine.
In other words, we have reduced to the case where $S_0$ is affine.
Next, we apply the induction principle again, this time to all
quasi-compact and quasi-separated opens of $X_0$, to reduce to the
case where $X_0$ is affine as well.

\medskip\noindent
If $X_0$ and $S_0$ are affine, the result follows from
More on Algebra, Lemma \ref{more-algebra-lemma-colimit-relatively-perfect}.
Namely, by Lemmas \ref{lemma-pseudo-coherent} and
\ref{lemma-affine-compare-bounded}
the statement is translated into computations of homs in the
derived categories of modules. Then
Lemma \ref{lemma-pseudo-coherent-affine}
shows that the complex of modules corresponding to $K_0$
is pseudo-coherent. And
Lemma \ref{lemma-tor-dimension-rel-affine}
shows that the complex of modules corresponding to $L_0$
has finite tor dimension over $\mathcal{O}_{S_0}(S_0)$.
Thus the assumptions of
More on Algebra, Lemma \ref{more-algebra-lemma-colimit-relatively-perfect}
are satisfied and we win.
\end{proof}

\end{document}
